\section{Packed Metadata}
\label{sec:packed}

\subsection{Organizing immutable records for reuse}

Once a source view is fixed, its metadata no longer needs to accommodate
in-place record updates. The producer can organize records for retrieval,
build their indexes once, and let readers reuse that organization for the
version's lifetime. Packed metadata applies this principle to directory
entries, common attributes, and file-data locators, placing related records
in bounded pages on object storage. It moves organization work to
construction while retaining demand-driven access to the resulting snapshot.

The snapshot describes a selected subtree and can directly back a standalone
read-only mount. It can also serve as the optional immutable lower of a
delta-based workspace, where the workspace mechanism resolves private
changes before consulting inherited records. This choice of representation
does not require a particular backend layout; native catalog records remain
an alternative lower representation.

A logical read group covers one parent directory and a range of raw entry
names in byte order. Its GroupMeta page co-locates each entry with common
attributes and data locators, so a name lookup can supply information needed
by subsequent attribute and data operations. Neighboring names may reuse
the same page. The raw page is bounded at 256 KiB, while a physical container
may hold several groups. This separates the unit of metadata organization
from the object that stores it. Range indexes route requests to these groups;
their height depends on page capacity and record count, rather than on
the source directory tree's depth.

\begin{figure}[t]
\centering
\begin{tikzpicture}[
  b/.style={draw,rounded corners=2pt,align=center,text width=6.5cm,minimum height=0.56cm,font=\small,fill=black!3},
  a/.style={-{Latex[length=1.7mm]},semithick}
]
\node[b] (root) {Pinned manifest and authenticated roots};
\node[b,below=0.35cm of root] (index) {Paged directory-name index\\$(\text{parent},\text{raw-name interval})$};
\node[b,below=0.35cm of index] (group) {GroupMeta page\\entry $+$ common attributes $+$ data locator};
\node[b,below=0.35cm of group] (frame) {Packed frame descriptor\\container offset, length, content verification};
\node[b,below=0.35cm of frame] (data) {Payload range read and verification};
\draw[a] (root) -- (index);
\draw[a] (index) -- (group);
\draw[a] (group) -- (frame);
\draw[a] (frame) -- (data);
\end{tikzpicture}

\caption{Representative packed-snapshot lookup/read path.
A name lookup supplies joint metadata and a data locator; reading referenced
data resolves authenticated descriptors before fetching and verifying payload ranges.}
\label{fig:packed}
\end{figure}

\subsection{Querying a pinned snapshot}

A reader pins a manifest identity and starts from its authenticated index
roots. For \code{lookup(parent,name)}, the directory-name index identifies
the group covering the key. The reader resolves its container reference and
fetches the GroupMeta byte range. It verifies the stored-byte digest before
decoding, then checks the decoded entry count and name boundaries against
the authenticated group reference. The matching record supplies the inode,
common attributes, and data locator (Figure~\ref{fig:packed}).

The query input determines which index is needed. Direct \code{getattr}
can obtain common attributes through an inode index, while enumeration
traverses ordered name ranges and uses counted indexes for directory-ordinal
navigation. A file read returns inline contents or follows a data locator
to an authenticated descriptor, fetches the referenced payload range, and
verifies it. Packed data frames define snapshot-level access and integrity
units; the backend serves these object requests using its own persisted
layout.

For supported snapshot queries, these paths obtain names, attributes, and
data locations without consulting per-file records in the mutable metadata
database. A writable workspace still performs its private upper checks and
control operations separately. Lower-query independence thus describes
where inherited records are obtained, while leaving workspace mutation and
lifetime coordination to their own protocols.

\subsection{Preserving identity and authenticated results}

Every returned record or payload must be connected by authenticated
references to the pinned manifest. Checking a fetched page's digest is one
part of this rule; checking its declared key boundaries and decoded contents
also prevents accepting the wrong records for a query. A missing referenced
object or failed authentication is an error, not evidence that a requested
name is absent.

The same pinned identity governs reuse. Authenticated pages and resolved
records remain valid for that version without tracking per-record updates.
Publishing a later version therefore does not change the records seen by
an existing reader. Caches can reuse these results, but their effectiveness
depends on locality: a cold lookup may traverse several index pages, and a
data read may require descriptor and payload requests beyond the metadata
page. Co-location provides an opportunity to reuse fetched information,
rather than a fixed number of GETs per operation.

\subsection{Publication, lifetime, and amortization}

Construction needs a consistent source view, established through a snapshot
or an explicit validation protocol. The producer groups records, uploads
their referenced objects, and produces the indexed manifest. Output
immutability alone cannot establish source consistency. Before publishing
a binding, the selected protocol must establish that its required objects
are durable and valid; producing a manifest and publishing its reference are
separate events. Readers retain their pinned versions, and reclamation must
preserve objects still required by retained versions or readers. Rebuilding
creates another version; ordinary file close does not rebuild a snapshot.

The access benefit depends on how much of each fetched page is useful.
Repeated nearby lookups can reuse joint metadata, whereas isolated lookups
may transfer records they never consume. Inlining file contents, compressing
pages, and co-packing payload regions also affect transfers, decoding, and
cache use. Their costs must be separated from the effect of metadata
co-location.

Construction, validation, publication, rebuilding, and reclamation all enter
the cost comparison. For one published version's lifecycle and $K$ uses,
we will measure each additive resource metric $c$ as
\begin{equation}
C_c(K)=C_{c,\mathrm{build}}+C_{c,\mathrm{publish}}
       +\sum_{i=1}^{K}C_{c,i}+C_{c,\mathrm{reclaim}}.
\end{equation}
The per-use terms include setup and access work for the chosen reuse unit;
construction and publication terms include their required validation.
CPU-seconds, transferred bytes, and monetary cost are separate metrics;
concurrent task wall-clock times are not summed into a batch makespan.
For repeated rebuilding, costs are accumulated over the resulting versions,
with shared work charged only once.
Reuse pays for construction only when the accumulated access savings exceed
these costs. Frequent rebuilding, weak name locality, or expensive cold
index walks can erase the benefit.

Immutable object-backed metadata already has strong
precedents~\cite{deltafs2021,cvmfs2011,nydus_rafs}. We will compare the joint
layout with immutable sorted-KV records at equal content and cache budgets,
and with compatible demand-loaded immutable-image systems. Standalone
read-only and workspace-lower use are separate comparisons. The sorted-KV
control tests layout organization and is not a reproduction of DeltaFS.
